\documentclass[a4paper, serif, 10pt]{moderncv}

%% moderncv
% style and color
\moderncvstyle{banking}
\moderncvcolor{black}

%% page layout/margins
\usepackage[top=2.5cm, bottom=2.5cm, left=1.5cm, right=1.5cm]{geometry}

\nopagenumbers{}

%% Babel config for Norwegian language
% ref:
% - https://latex3.github.io/babel/guides/locale-norwegian.html
\usepackage[norsk]{babel}

%% fontspec
\usepackage{fontspec}

\IfFontExistsTF{Linux Libertine}{
  \setmainfont[Ligatures={TeX, Common}, Numbers=OldStyle]{Linux Libertine}
}{}
\IfFontExistsTF{Linux Libertine O}{
  \setmainfont[Ligatures={TeX, Common}, Numbers=OldStyle]{Linux Libertine O}
}{}

%% CTeX
% CJK support, used to show CJK characters in the resume
%
% - fontset=none: disable builtin fontset but instead set the CJK font manually
% - heading=false: disable ctex heading
% - punct=kaiming: use kaiming punctuations styles for CJK
% - scheme=plain: use plain scheme, do not override \normalsize font size
% - space=auto: space settings for CJK characters
%
% ref:
% - http://ctan.mirrorcatalogs.com/language/chinese/ctex/ctex.pdf
\usepackage[UTF8, heading=false, punct=kaiming, scheme=plain, space=auto]{ctex}

\IfFontExistsTF{Noto Serif CJK SC}{
  \setCJKmainfont{Noto Serif CJK SC}
}{}
\IfFontExistsTF{Noto Sans CJK SC}{
  \setCJKsansfont{Noto Sans CJK SC}
}{}

%% line spacing
\usepackage{setspace}
\setstretch{1.125}

%% URL styling - use normal text instead of monospace
\urlstyle{same}

% Auto-underline all links
% Defer until \begin{document} so hyperref is already loaded
\AtBeginDocument{%
  \let\oldhref\href
  \renewcommand{\href}[2]{\oldhref{#1}{\underline{#2}}}%
}

%% Patch \httplink and \httpslink to handle full URLs
% The original moderncv macros always prepend http:// or https:// to URLs.
% This causes issues when the URL already has a protocol (e.g., https://example.com)
% resulting in malformed URLs like https://https://example.com.
% This patch checks if the URL already contains "://" and uses it directly if so.
% Note: We use \str_set:Nx (x-expansion) to expand macros like \@homepage first.
\ExplSyntaxOn
\str_new:N \g_yamlresume_url_str

\RenewDocumentCommand{\httpslink}{O{}m}{
  \str_set:Nx \g_yamlresume_url_str {#2}
  \str_if_in:NnTF \g_yamlresume_url_str {://}
    { \url{#2} }
    { \url{https://#2} }
}

\RenewDocumentCommand{\httplink}{O{}m}{
  \str_set:Nx \g_yamlresume_url_str {#2}
  \str_if_in:NnTF \g_yamlresume_url_str {://}
    { \url{#2} }
    { \url{http://#2} }
}
\ExplSyntaxOff

\name{Ingrid Solberg}{}
\title{Kundesuksessansvarlig}
\phone[mobile]{+47 22 33 44 55}
\email{ingrid.solberg@example.com}
\homepage{https://www.ingridsolberg.no}

\address{Storgata 45, Oslo, Oslo, Norge, 0182}{}{}

\extrainfo{{\small \faLinkedin}\ \href{https://www.linkedin.com/in/ingridsolberg}{@ingridsolberg} {} {} {} • {} {} {} 
{\small \faGithub}\ \href{https://github.com/ingridsolberg}{@ingridsolberg}}

\begin{document}

\maketitle

\section{Grunnleggende informasjon}

\cvline{}{Resultatorientert kundesuksessansvarlig med over 6 års erfaring fra SaaS- og teknologibransjen. Jeg brenner for å bygge sterke kundeforhold, øke kundetilfredshet og redusere churn gjennom proaktiv oppfølging og datadrevne innsikter.
%
\begin{itemize}
\item Dokumentert suksess med å øke netto opprettholdelse (NRR) og kundelivstidsverdi (LTV)
\item Sterk kompetanse innen onboarding, opplæring og strategisk rådgivning
\item Erfaren med CRM-verktøy som Salesforce, HubSpot og Gainsight
\item Flytende i norsk og engelsk, med gode kommunikasjons- og presentasjonsevner
\end{itemize}}
\section{Utdanning}

\cventry{aug. 2016–juni 2018}
        {Master, Økonomi og administrasjon, Poeng: 4.6}
        {Universitetet i Oslo}
        {\url{https://www.uio.no}}
        {}
        {\begin{itemize}
\item Spesialisering innen strategi, markedsføring og organisasjonsutvikling
\item Gjennomførte flere prosjekter med fokus på kundeinnsikt og forretningsvekst
\item Utviklet sterke analytiske ferdigheter og evne til å formidle komplekse sammenhenger
\end{itemize}
\textbf{Kurs}: Strategisk markedsføring, Organisasjonspsykologi, Forretningsutvikling, Dataanalyse for beslutningstaking, Prosjektledelse, Kundeadferd og relasjonsledelse}

\cventry{aug. 2013–juni 2016}
        {Bachelor, Markedsføring og merkevareledelse, Poeng: 4.2}
        {Handelshøyskolen BI}
        {\url{https://www.bi.no}}
        {}
        {\begin{itemize}
\item Grunnleggende forståelse for markedsføring, salg og kundeadferd
\item Deltok i studentorganisasjoner og arrangerte karrieredager
\end{itemize}
\textbf{Kurs}: Markedsføringsledelse, Forbrukeratferd, Salgsledelse, Statistikk, Kommunikasjon og presentasjonsteknikk}
\section{Arbeidserfaring}

\cventry{mars 2021–Nå}
        {Kundesuksessansvarlig}
        {Visma}
        {\url{https://www.visma.no}}
        {}
        {\begin{itemize}
\item Ansvarlig for en portefølje på over 40 B2B-kunder innen SaaS-løsninger for HR og lønn
\item Økte netto opprettholdelse (NRR) fra 92 \% til 108 \% gjennom proaktiv rådgivning og verdiskapende oppfølging
\item Ledet onboarding av nye kunder, inkludert opplæring, implementering og milepælsplaner
\item Reduserte churn med 25 \% ved å etablere tidlige varslingssystemer og regelmessige helsesjekker
\item Fungerte som bindeledd mellom kunder og produktteam for å sikre kundedrevet produktutvikling
\end{itemize}
\textbf{Nøkkelord}: Kundeoppfølging, Onboarding, SaaS, Churn, Salesforce}

\cventry{jan. 2019–feb. 2021}
        {Kundesuksessspesialist}
        {Kahoot!}
        {\url{https://kahoot.com}}
        {}
        {\begin{itemize}
\item Håndterte kundehenvendelser og support for over 200 skoler og bedrifter i Norden
\item Utviklet og gjennomførte opplæringswebinarer som økte produktadopsjonen med 35 \%
\item Bidro til å bygge opp en kundesuksessavdeling fra grunnen av, med fokus på skalerbare prosesser
\item Brukte HubSpot til å spore kundeinteraksjoner og identifisere oppsalgspotensial
\end{itemize}
\textbf{Nøkkelord}: Kundestøtte, Webinarer, HubSpot, Adopsjon, Norden}

\cventry{juni 2017–jan. 2019}
        {Kundebehandler}
        {Telenor}
        {\url{https://www.telenor.no}}
        {}
        {\begin{itemize}
\item Betjente bedriftskunder i telefon og e-post, med fokus på rask og profesjonell problemløsning
\item Oppnådde en kundetilfredshetscore (CSAT) på 4,8 av 5 i snitt
\item Identifiserte behov for tilleggstjenester og videreformidlet salgsmuligheter til selgerteamet
\end{itemize}
\textbf{Nøkkelord}: Kundeservice, CSAT, Bedriftsmarked, Problemløsning}
\section{Språk}

\cvline{Norsk}{Morsmål eller tospråklig nivå \hfill \textbf{Nøkkelord}: Morsmål}
\cvline{Engelsk}{Fullt profesjonelt nivå \hfill \textbf{Nøkkelord}: Forretningsflyt, TOEFL 105}
\section{Ferdigheter}

\cvline{Kundesuksess}{Ekspert \hfill \textbf{Nøkkelord}: Onboarding, Kundereiser, Churn-analyse, Oppsalg, Helseundersøkelser}
\cvline{CRM og verktøy}{Avansert \hfill \textbf{Nøkkelord}: Salesforce, HubSpot, Gainsight, Zendesk, Looker}
\cvline{Kommunikasjon}{Ekspert \hfill \textbf{Nøkkelord}: Presentasjoner, Forhandlinger, Workshop, Relasjonsbygging}
\cvline{Dataanalyse}{Avansert \hfill \textbf{Nøkkelord}: Excel, SQL, Dashboard, KPI-oppfølging, Segmentering}
\section{Utmerkelser}

\cventry{des. 2023}
        {Visma}
        {Årets Kundesuksessansatt}
        {}
        {}
        {Tildelt for eksepsjonell innsats med å øke kundetilfredshet og opprettholdelse i et krevende marked.}
\section{Sertifikater}

\cventry{mai 2022}
        {SuccessCoaching}
        {Certified Customer Success Manager (CCSM)}
        {\url{https://www.successcoaching.com/certification}}
        {}
        {}

\cventry{okt. 2021}
        {Salesforce}
        {Salesforce Certified Administrator}
        {\url{https://trailhead.salesforce.com}}
        {}
        {}
\section{Publikasjoner}

\cventry{apr. 2023}
        {Hvordan bruke kundedata til å drive lønnsom vekst}
        {Magma – Tidsskrift for økonomi og ledelse}
        {\url{https://www.magma.no}}
        {}
        {\begin{itemize}
\item Artikkelen utforsker hvordan SaaS-selskaper kan bruke kundedata for å forbedre opprettholdelse og øke omsetning.
\item Presenterer en modell for proaktiv kundeoppfølging basert på helsescore og bruksmønstre.
\end{itemize}}
\section{Prosjekter}

\cventry{sep. 2022–feb. 2023}
        {Utviklet et dashboard for å visualisere kundehelse og identifisere churn-risiko i sanntid.}
        {Kundehelse-dashboard}
        {\url{https://www.ingridsolberg.no/prosjekter/kundehelse}}
        {}
        {\begin{itemize}
\item Integrerte data fra Salesforce og produktbruk for å lage en helhetlig kundehelsescore
\item Automatiserte varsler til kundesuksessansvarlige ved fall i engasjement
\item Reduserte tiden brukt på manuell rapportering med 60 \%
\end{itemize}
\textbf{Nøkkelord}: Dashboard, Kundehelse, Automatisering}

\cventry{juni 2021–jan. 2022}
        {Redesignet onboardingprogrammet for nye kunder for å øke tid-til-verdi.}
        {Onboarding 2.0}
        {\url{https://www.ingridsolberg.no/prosjekter/onboarding}}
        {}
        {\begin{itemize}
\item Kartla kundereisen og identifiserte flaskehalser i implementeringsfasen
\item Utviklet nye maler, sjekklister og opplæringsvideoer
\item Kortet ned gjennomsnittlig onboardingtid fra 45 til 28 dager
\end{itemize}
\textbf{Nøkkelord}: Onboarding, Prosessforbedring, Tid-til-verdi}
\section{Interesser}

\cvline{Friluftsliv}{Fjellturer, Ski, Kajakk}
\cvline{Samfunn}{Mentorskap, Bærekraft, Teknologi}
\section{Frivillig arbeid}

\cventry{sep. 2020–Nå}
        {Mentor}
        {Ungt Entreprenørskap Norge}
        {\url{https://www.ue.no}}
        {}
        {\begin{itemize}
\item Veiledet unge gründere i utvikling av forretningsideer og kundeforståelse
\item Holdt workshops i kundesegmentering og verdivurdering
\item Bidro til å styrke studenters presentasjons- og samarbeidsevner
\end{itemize}}
    
\end{document}